\section{研究背景与课题意义}

\subsection{研究背景}

随着互联网应用和数字娱乐平台的快速发展，电影、短视频、音乐、新闻和商品等信息资源呈现爆炸式增长。
用户在获得更丰富选择的同时，也面临越来越明显的信息过载问题。以在线电影平台为例，平台中通常包含
大量不同年代、类型、国家和风格的影片，如果只依靠用户主动搜索、人工分类或热门排行榜，用户很难在
有限时间内快速找到真正符合个人兴趣的内容。与此同时，平台积累了大量用户评分、浏览、点击和收藏等
行为数据，这些数据中包含了用户偏好和群体兴趣之间的关联。如果不能对这些数据进行有效建模和利用，
就会造成数据价值的浪费。推荐系统研究表明，用户--物品交互数据能够有效支持个性化信息过滤与内容发现
\cite{adomavicius2005toward}。

推荐系统正是在这样的背景下产生并不断发展的。推荐系统通过分析用户历史行为、物品属性以及用户与物品
之间的交互关系，从大量候选内容中筛选出更可能满足用户需求的结果。与传统搜索系统相比，推荐系统不
要求用户明确输入完整需求，而是能够根据历史偏好主动提供个性化内容；与简单排行榜相比，推荐系统又
能够考虑不同用户之间的兴趣差异，从而提高内容发现效率。因此，推荐系统已经成为现代 Web 应用、电子
商务、在线视频平台和数据智能系统中的重要基础能力。工业界典型系统如 Amazon 的 item-to-item
协同过滤也说明，相似物品推荐可以在大规模在线服务中取得较好的实用效果\cite{linden2003amazon}。

电影推荐是推荐系统研究和教学中非常典型的应用场景。一方面，电影具有标题、类型、评分、标签等较清晰
的数据结构，便于构建可解释的推荐流程；另一方面，用户对电影的评分和观看行为能够较好反映兴趣偏好，
适合用于协同过滤、相似度计算、矩阵分解和离线评估等实验。MovieLens 数据集作为推荐系统领域常用的
公开数据集，提供了稳定的电影评分数据和较规范的数据格式，既适合算法实验，也适合作为课程设计的数据
基础\cite{harper2015movielens}。在算法方面，基于物品的协同过滤、矩阵分解、隐式反馈建模和排序优化等
方法均是电影推荐实验中常见的技术路线\cite{sarwar2001item,koren2009matrix,hu2008collaborative,rendle2009bpr}。

在实际应用中，推荐系统不仅是一个算法问题，也是一个完整的软件工程问题。一个可用的推荐系统通常需要
完成数据导入、数据存储、推荐计算、接口服务、缓存加速、前端展示、指标评估和部署运维等多个环节。
如果只在 Notebook 中计算推荐结果，虽然可以验证算法公式，但难以体现推荐系统在线服务的真实流程。
因此，本课程设计选择将电影推荐算法与 Web 服务、数据库、缓存和容器化部署结合起来，构建一个能够
运行、交互和评估的云端电影推荐系统原型。推荐系统评估研究也指出，除推荐准确率外，排序质量、响应性能
和实验可复现性同样会影响系统效果判断\cite{herlocker2004evaluating}。

本课程设计以 FastAPI Movie Recommender 项目为基础，将原有单体 Demo 扩展为一个基于 Docker Compose
部署的电影推荐系统。系统采用 FastAPI 提供后端接口，使用 Streamlit 展示交互页面，使用 PostgreSQL
保存 MovieLens 数据，使用 Redis 缓存推荐结果，并提供功能测试和抽样离线评估能力。该设计既关注推荐
算法的基本效果，也关注系统在真实 Web 服务中的可部署性、可维护性和可扩展性。

\subsection{课题意义}

本课题围绕电影推荐系统展开，既具有一定的应用价值，也具有明确的课程实践意义。通过该系统，可以把
推荐算法、后端开发、数据库设计、缓存机制、前端展示和容器化部署联系起来，形成一个相对完整的工程
闭环。系统能够根据用户选择的电影返回相似电影，为电影内容检索、兴趣发现和个性化推荐提供辅助；
同时，推荐结果中包含相似度得分和推荐原因，便于观察协同过滤思想在具体场景中的应用过程。

从课程设计角度看，本课题的价值不只在于实现一个推荐算法，而在于把算法放入真实系统环境中进行验证。
MovieLens 数据集提供了用户--电影评分交互数据，PostgreSQL 和 Redis 分别承担数据持久化与缓存加速，
FastAPI 和 Streamlit 分别承担接口服务与页面展示，Docker Compose 则负责多服务统一编排。通过这一
过程，可以更完整地理解推荐系统从数据处理、推荐计算到在线展示和实验评估的实现链路。

\subsection{课程设计目标}

本课程设计的总体目标是完成一个可以运行、可以交互、可以部署、可以初步评估的电影推荐系统原型。
系统不仅需要能够给出推荐结果，还需要具备清晰的数据链路、接口设计、页面展示和实验评价方法。
围绕这一总体目标，本文主要完成以下任务：

\begin{enumerate}[leftmargin=2em]
  \item 完成 MovieLens 数据集的导入与管理，建立电影、用户、评分、点击事件和实验结果等基础数据表，为推荐计算和结果分析提供数据支撑。
  \item 设计并实现电影检索、相似电影推荐、推荐解释和 Demo 点击统计等核心功能，使用户能够在前端页面中完成完整的推荐交互流程。
  \item 使用 FastAPI、Streamlit、PostgreSQL、Redis 和 Docker Compose 搭建可部署的推荐系统，实现后端接口、前端展示、数据存储、缓存加速和多容器编排。
  \item 设计功能测试与抽样离线评估方案，从推荐命中情况、排序效果和响应延迟等角度分析系统表现，并总结后续可改进的方向。
\end{enumerate}

\newpage
